
\subsection{Primitive large sieve and the powers of $e$}

We use the classical primitive multiplicative large sieve in the form
\cite{IwaniecKowalski2004,MontgomeryVaughan2007}
\begin{equation}\label{eq:primitive-ls}
 \sum_{f\le Q}\frac f{\phi(f)}
 \sum_{\chi^*\, (\mathrm{mod}\ f)}^{\mathrm{prim}}
 \left|\sum_m c_m\chi^*(m)\right|^2
 \le (Q^2+N)\sum_m|c_m|^2,
\end{equation}
when $(c_m)$ is supported in an interval of length at most $N$.  Since
$f/\phi(f)\ge1$, the weight may be discarded for an upper bound.

For fixed $e$, take
\[
 c_m=a_{em}\1_{D/e<m\le y/e}.
\]
Its support has length at most $2D/e$, and
\begin{equation}\label{eq:c-energy-rough}
                  \sum_m|c_m|^2\le\frac{2D}{e}
\end{equation}
whenever it is nonempty.  Applying \cref{eq:primitive-ls} with $Q=2F$ to
\cref{eq:VF}, we obtain
\begin{align}
 V_F(y)
 &\ll 4^j\sum_{e\ge1}\frac1{Fe}
       \left(4F^2+\frac{2D}{e}\right)\frac{2D}{e}\notag\\
 &\ll 4^j\left(
       FD\sum_{e\ge1}e^{-2}
       +\frac{D^2}{F}\sum_{e\ge1}e^{-3}
       \right)\notag\\
 &\ll_j FD+\frac{D^2}{F}.
 \label{eq:VF-final}
\end{align}
The powers $e^{-2}$ and $e^{-3}$ are load-bearing: one $e^{-1}$ comes from
complementary-modulus multiplicity, and the remaining powers come from support
length and coefficient energy.

\subsection{Dyadic conductor summation}

By \cref{eq:f-lower}, the dyadic conductors satisfy $P\le F<2R$.  Hence
\begin{equation}\label{eq:F-sums}
       \sum_F FD\ll DR,
 \qquad
       \sum_F\frac{D^2}{F}\ll\frac{D^2}{P}.
\end{equation}
Both are geometric sums; neither costs a logarithm.  Summing
\cref{eq:VF-final} proves \cref{eq:roughvariance}.

Tracking the factors in
\cref{eq:phi-lower,eq:imprimitive-cauchy,eq:r-count} and using
$\sum e^{-2}<2$, $\sum e^{-3}<2$, together with the standard leading constant
in \cref{eq:primitive-ls}, shows that $2^{2j+8}$ is a deliberately generous
coefficient.  The theorem itself is recorded with $\ll_j$ because endpoint
normalisations of the large sieve vary across references.
\qed

\begin{remark}[What is and is not standard]
The primitive large sieve is standard.  The derivation above, including the
roughness gain $D^2/P$, the exact imprimitive correction, and its insertion into
the physical common-sequence dictionary, is the new derivation in this
programme.  No claim is made that the abstract consequence
\cref{eq:roughvariance} is absent from the literature.
\end{remark}

\section{Physical transference of the variance theorem}\label{sec:transference}

Fix a cell and write $i=(q,k)$, with $q$ in the physical five-prime modulus set
and $k\asymp K$.  Since $(q,k)=1$, the determinant congruence selects the
reduced residue
\begin{equation}\label{eq:ai}
                        a_i\equiv2\bar k\pmod q.
\end{equation}
For fixed $q$ and fixed reduced residue $a$, the number of $k$ in a dyadic
interval of length $K$ satisfying \cref{eq:ai} is
\begin{equation}\label{eq:residue-multiplicity}
                        O(1+K/R).
\end{equation}
Indeed, $k$ is fixed modulo $q$ and $q\asymp R$.

Zero-extend the centred target to a cyclic group of order $Q_i\asymp D/q$,
with $Q_i$ at least twice the line length, and let $c_i(\xi)$ be as in
\cref{eq:cxi}.  The physical line coordinate $t_i(d)$ is affine with slope
$1/q$; since $qQ_i\asymp D$,
\begin{equation}\label{eq:phase-derivative}
 \left|\frac{d}{dy}\eord{Q_i}{\xi t_i(y)}\right|
                       \ll\frac{|\xi|}{D}.
\end{equation}
Only this Fourier atom is differentiated.  The six-prime coefficient is not
assumed differentiable as a function of $t$.

For $1\le H<Q_i/2$, define the exact band-limited first moment
\begin{equation}\label{eq:LleH}
 L_i^{\le H}=
 \sum_{\substack{d\ \mathrm{in\ the\ physical\ dyadic\ line}}}
 \mu(d)\sum_{0<|\xi|\le H}c_i(\xi)
                   \eord{Q_i}{\xi t_i(d)}.
\end{equation}

\subsection{The nonprincipal progression part}

Let $M_i^\circ(\xi)$ denote the sum in $d$ at frequency $\xi$ after the
principal progression mean is subtracted.  Stieltjes summation gives
\begin{equation}\label{eq:stieltjes}
 M_i^\circ(\xi)
   =\int_D^{2D}\eord{Q_i}{\xi t_i(y)}\,dE_q(a_i;y),
\end{equation}
with a bounded number of analogous endpoint terms when the physical interval is
cut into subintervals.  Minkowski's inequality in the $\ell^2(i)$ norm,
\cref{eq:phase-derivative}, the residue multiplicity
\cref{eq:residue-multiplicity}, and \cref{thm:roughvariance} imply
\begin{equation}\label{eq:Mxi-l2}
 \sum_i|M_i^\circ(\xi)|^2
 \ll_j (1+K/R)(1+|\xi|)^2
             \left(DR+\frac{D^2}{P}\right).
\end{equation}
To see the factor $1+|\xi|$, the endpoint term costs one copy of
$\sup_y\norm{E(y)}_2$, while the integral costs
$D\cdot(|\xi|/D)\sup_y\norm{E(y)}_2$.

For each $i$, Cauchy--Schwarz in frequency and \cref{eq:c-energy} give
\[
 |L_i^{\le H,\circ}|^2
 \le B_0^2\sum_{0<|\xi|\le H}|M_i^\circ(\xi)|^2.
\]
There are $O(RK)$ pairs $(q,k)$.  Summing in $i$ and applying Cauchy--Schwarz
once more gives
\begin{align}
 \sum_i|L_i^{\le H,\circ}|
 &\ll_j B_0\sqrt{RK}\sqrt{1+K/R}
 \left(
   \sum_{0<|\xi|\le H}(1+|\xi|)^2
   \left(DR+\frac{D^2}{P}\right)
 \right)^{1/2}\notag\\
 &\ll_j B_0(1+H)^{3/2}\sqrt{RK}\sqrt{1+K/R}
          \sqrt{DR+\frac{D^2}{P}}.
 \label{eq:nonprincipal-transfer}
\end{align}
For $K\le R$, the multiplicity factor is absolute.

\subsection{The principal progression part}

The principal term cannot be discarded: it is introduced and subtracted by an
exact character expansion.  For every fixed $J>0$, the classical zero-free
region for $\zeta(s)$ gives the uniform Mertens estimate
\begin{equation}\label{eq:mertens}
             \sup_{D\le y\le2D}
             \left|\sum_{D<n\le y}\mu(n)\right|
             \ll_J D\LL^{-J}
\end{equation}
for the polynomial-sized $D$ occurring here; see, for example,
\cite{MontgomeryVaughan2007}.  Now use the exact inclusion--exclusion identity
\[
             \1_{(n,q)=1}=\sum_{e\mid(n,q)}\mu(e).
\]
The $e=1$ term is handled by \cref{eq:mertens}.  Every nontrivial divisor
$e\mid q$ is at least $P$, and there are at most $2^j$ divisors.  Trivial
counting therefore gives the corrected bound
\begin{equation}\label{eq:principal-coprime}
 \sup_{D\le y\le2D}
 \left|\sum_{\substack{D<n\le y\\(n,q)=1}}\mu(n)\right|
 \ll_{J,j}D\LL^{-J}+O_j\!\left(2^j\left(\frac DP+1\right)\right).
\end{equation}
The first draft's literal coefficient $j(D/P+1)$ is not justified by direct
deletion because intersections of prime multiples must be included.  The
$2^j$ inclusion--exclusion cost is the audited correction; since $j=5$ in the
physical source, it changes no exponent.

A second Stieltjes summation and \cref{eq:phase-derivative} show that, for each
$i$,
\begin{equation}\label{eq:principal-Mxi}
 |M_i^{\mathrm{prin}}(\xi)|
 \ll_{J,j}(1+|\xi|)\frac DR
       \left(\LL^{-J}+P^{-1}+D^{-1}\right).
\end{equation}
Using \cref{eq:c-energy}, summing over the $O(RK)$ indices, and noting
$\sum_{|\xi|\le H}(1+|\xi|)^2\ll(1+H)^3$, we obtain
\begin{equation}\label{eq:principal-transfer}
 \sum_i|L_i^{\le H,\mathrm{prin}}|
 \ll_{J,j}B_0DK(1+H)^{3/2}
       \left(\LL^{-J}+P^{-1}+D^{-1}\right).
\end{equation}
The same calculation allows the line-constant target frequency $\xi=0$.

If the source sequence is $\mu(n)\1_{2\nmid n}$, write
\[
 M_{\mathrm{odd}}(x)=\sum_{\substack{n\le x\\2\nmid n}}\mu(n).
\]
The exact relation
\[
       M(x)=M_{\mathrm{odd}}(x)-M_{\mathrm{odd}}(x/2)
\]
gives
\begin{equation}\label{eq:odd-recursion}
       M_{\mathrm{odd}}(x)=\sum_{\ell\ge0}M(x/2^\ell),
\end{equation}
with a finite sum.  A conservative uniform implementation pays one additional
factor $O(\log D)$; equivalently, one decreases the available Mertens exponent
by one.  This is the source of the corrected global log ledger in
\cref{sec:auditconstants}.

\subsection{Normalisation}

Combining \cref{eq:nonprincipal-transfer,eq:principal-transfer} and using
$DK\asymp X$ gives
\begin{align}
 \frac{\sqrt{RK}\sqrt{DR+D^2/P}}{DK}
 &=\left(\frac{R^2}{DK}+\frac{R}{KP}\right)^{1/2}\notag\\
 &=\left(\frac{R^2}{X}+\frac{R}{KP}\right)^{1/2}.
 \label{eq:normalised-deficit}
\end{align}
Thus \cref{thm:transference-main} follows.  For $K>R$, the squared
nonprincipal deficit is
\begin{align}
 \left(1+\frac KR\right)
 \left(\frac{R^2}{X}+\frac{R}{KP}\right)
  &=\frac{R^2}{X}+\frac RD+\frac{R}{KP}+\frac1P,
 \label{eq:four-term-deficit}
\end{align}
where $DK=X$ was used.

\section{Physical closure of the former Sector A}\label{sec:sectorA}

Set $P=P_*$.  In the range
\begin{equation}\label{eq:enlarged-region}
             K\ge\frac{R}{\sqrt{P_*}},\qquad K\le R,
\end{equation}
we have, by $R\ll X^{5/11}$,
\begin{equation}\label{eq:sector-deficits}
       \frac{R^2}{X}\ll X^{-1/11+o(1)},
       \qquad
       \frac{R}{KP_*}\le P_*^{-1/2}.
\end{equation}
Therefore the square root in \cref{eq:transference-main} is
\begin{equation}\label{eq:delta}
 \ll X^{-1/22+o(1)}+P_*^{-1/4}
 \ll X^{-\delta+o(1)}
\end{equation}
for every fixed
\begin{equation}\label{eq:delta-range}
        0<\delta<\min\left(\frac1{22},\frac{\epsstar}{2}\right).
\end{equation}
If $H\le\LL^b$, the factor $(1+H)^{3/2}$ is at most a fixed power of
$\LL$, and the power saving in \cref{eq:delta} absorbs every declared finite
logarithmic cost.  The principal term is made arbitrarily small by choosing
$J$ sufficiently large.  This proves \cref{cor:sectorA}.

For cells with $K>R$, the extra term in \cref{eq:four-term-deficit} obeys
\[
       \frac RD\ll X^{-1/11+\epssafe}
\]
under \cref{eq:lowd-cutoff}.  If $\epssafe<1/11$, the overlap remains
power-saving, with any
\begin{equation}\label{eq:delta-Kgreater}
 0<\delta<\min\left(
       \frac{1/11-\epssafe}{2},\frac{\epsstar}{2}
       \right).
\end{equation}
This extension is included only on cells actually permitted by the inherited
